\documentclass[12pt,aspectratio=169]{beamer}

%%%
% Einbinden des UniWue Theme 
% Weitere Hinweise unter templates/styles/beamerUniWue/README
%%%
\usetheme{UniWue}
\usefonttheme{professionalfonts}


\usepackage[utf8]{inputenc}     % Umlaute und ähnliches möglich
\usepackage[english]{babel}		% Deutsche Rechtschreibung
\usepackage[T1]{fontenc}        % verbessert Worttrennung (westeuropäischer Codierung)
\usepackage{lmodern}			% Schrift wirkt (in pdf-Ausgabe) fließender
\usepackage{amsmath}			% Matheformeln
\usepackage{amsfonts}           % zusätzliche mathematische Symbole 
\usepackage{amssymb}            % zusätzliche mathematische Symbole 
\usepackage{graphicx}			% Bilder einfügen
\usepackage[version=4]{mhchem} 	% diverse Chemie-Funktionen
\usepackage{pdfpages} 			% macht PDFs importierbar
\usepackage{floatflt,epsfig}    % für die Platzierung von Abbildungen etc
\usepackage{svg}                % SVG einbinden
\usepackage{lipsum}
\usepackage{subcaption}				% damit Bilder nebeneinander sein können
\usepackage{esvect}					% Vektorpfeil
\usepackage{hyperref}				% für \url -> Internetseiten
\usepackage{appendixnumberbeamer}	% Gesamtseitenzahlen ab \appendix nicht erhöht
\usepackage{csquotes}				% für Bib
\usepackage[style=numeric, sorting=nty, maxcitenames=1, bibwarn=true, giveninits=true, isbn=false, url=true, doi=false, block=none]{biblatex}	% für die Quellen über bibtex
\usepackage{epstopdf}               % eps,gif,png zu pdf - siehe epstopdfDeclareGraphicsRule
%\usepackage{subcaption}             % für Subfigure
\usepackage{xcolor}                  % für Farben
\usepackage{tabulary}
\usepackage{caption}                  % um mit * Abb. punkt zu entfernen


\usepackage{siunitx}                %SI-Funktion
  \sisetup{
  range-phrase = {-},
  output-decimal-marker = {.}, 
  group-separator = {}, 
  exponent-product = \ensuremath{{}\cdot{}}, 
  per-mode = reciprocal, % mit fraction kann Bruchschreibweise erhalten werden, mit reciprocal negativer Exponent
  inter-unit-product = \ensuremath{{}\cdot{}},
  range-units = single,
  mode = math,
  detect-all,  
  number-unit-separator=\text{\,},
%   math-rm=\mathsf,
%   text-rm=\sffamily,
  }%
  
  
\makeatletter
\newcommand{\showfontsize}{\f@size{} pt}
\makeatother

\newcommand*\mycommand[1]{\texttt{\emph{#1}}}
\newcommand{\itemcolor}[1]{% Update list item colour
  \renewcommand{\makelabel}[1]{\color{#1}\hfil ##1}}


%%% Damit auch bei \mathrm keine Serifenschriftart verwendet wird:
\let\lmodernsfdefault\sfdefault
\let\rmdefault\sfdefault
\let\sfdefault\lmodernsfdefault
\renewcommand*{\familydefault}{\rmdefault}

%%% Veränderungen zum Zitieren in der Fußnote:
\newcommand*{\quelle}[1]{\par\raggedleft\tiny Quelle:~#1}
\renewcommand{\footnotesize}{\tiny}
\renewcommand{\thefootnote}{} %keine Zahl für Fußnote
\renewcommand{\footnoterule}{} %Löscht Linie über Fußnote
\renewcommand{\footnotesep}{-5pt}
% Damit alles in der Fußnote linksbündig ist:
%\usepackage{libertine} 
%\usepackage{microtype}
\makeatletter
\renewcommand\@makefntext[1]{%
    \parindent 1em%
    \noindent\hbox{\@thefnmark}\hspace*{0em plus 0pt minus 0pt}#1}
\makeatother

\newcommand{\myrule}[3]{\textcolor{#1}{\rule{#2}{#3}}}
\newcommand{\ca}{\raisebox{-0.75ex}{\~{}}} % für ~ im Text

\addbibresource{beamerbibliography.bib}

\setbeamertemplate{navigation symbols}{}

\author{Autor}
\title{Vorlage für Präsentationen der JMU in LaTeX}
%\setbeamercovered{transparent} 
%\logo{} 
\institute{Julius-Maximilians-Universität Würzburg} 
\date{Date} 
\subject{Subject} 
\titlegraphic{\includegraphics[width=\textwidth]{Bild.jpg}}

\begin{document}


{
\setbeamertemplate{footline}{} 
\begin{frame}
    \vspace{-1.5cm} % Hier anpassen der Höhe der Titlepage auf der Folie
    \titlepage
\end{frame}
}

\section*{Outline}
\begin{frame}
  \frametitle{Outline}
  \tableofcontents[hidesubsections]
\end{frame}

\AtBeginSection[]
{
  \frame<handout:0>
  {
    \frametitle{Outline}
    \tableofcontents[currentsection,hideallsubsections]
  }
}

\AtBeginSubsection[]
{
  \frame<handout:0>
  {
    \frametitle{Outline}
    \tableofcontents[sectionstyle=show/hide,subsectionstyle=show/shaded/hide]
  }
}

\newcommand<>{\highlighton}[1]{%
  \alt#2{\structure{#1}}{{#1}}
}

\newcommand{\icon}[1]{\pgfimage[height=1em]{#1}}



%%%%%%%%%%%%%%%%%%%%%%%%%%%%%%%%%%%%%%%%%
%%%%%%%%%% Content starts here %%%%%%%%%%
%%%%%%%%%%%%%%%%%%%%%%%%%%%%%%%%%%%%%%%%%



\section{Introduction}

\begin{frame}
  \frametitle{Prerequisites \& Goals}
  \framesubtitle{Knowledge}
  \begin{block}{LaTeX}
  \begin{itemize}
    \item Obviously some basic LaTeX knowledge is necessary
    \item Some more features will be provided here
  \end{itemize}
  \end{block}

  \begin{block}{Beamer}
  \begin{itemize}
    \item You'll learn them by looking at this presentation source
  \end{itemize}
  \end{block}

  \begin{block}{Goal}
  \begin{itemize}
    \item Learn how to make well structured slides
    \item Using a beautiful theme (congrats to the Oxygen team!)
    \item Take over the world
    \item Relax...
  \end{itemize}
  \end{block}
\end{frame}

\section{Basic structuring}
\begin{frame}
  \frametitle{Sections, Frames and Blocks}
  \framesubtitle{Put everything into boxes}

  The current section is "Basic structuring". 

  \begin{block}{A beautiful block}
  A block has a title, and some content. You can put in a block
  almost everything you want that is provided by LaTeX. For example
  math works as usual:
    \begin{equation}
    \sum_{i=1}^n i = \frac{n \times (n+1)}{2}
    \end{equation}
  \end{block}

  Also works outside a block:
  \begin{equation}
  \sum_{i=1}^n i^2 = \frac{n \times (n+1) \times (2n+1)}{6}
  \end{equation}
\end{frame}

\begin{frame}
  \frametitle{Different type of blocks}
  \framesubtitle{Weeeee! Colors!!}
  \begin{block}{Standard block}
  \begin{itemize}
    \item A standard block, used for grouping
    \item Obviously can contain itemizes too...
    \begin{itemize}
      \item And nested itemizes...
      \item of course!
    \end{itemize}
  \end{itemize}
  \end{block}
  \begin{alertblock}{Alert block}
  WARNING: Something very important inside this block!
  \end{alertblock}
  \begin{example}
  Note that examples are displayed as a special block...
  \end{example}
\end{frame}

\section{Fancy features}
\begin{frame}
  \frametitle{Highlighting}
  \framesubtitle{Hey! Look here! here!}

  \begin{block}{A regular block}
  \begin{itemize}
    \item Normal text
    \item \highlighton{Highlighted text} to draw attention
    \item \alert{"Alert'ed" text} to spot very important information
    \item Alternatively you can
    \begin{itemize}
      \alert{\item "Alert" the item itself}
      \highlighton{\item Or "Highlight" it}
    \end{itemize}
  \end{itemize}
  \end{block}
  \begin{alertblock}{If it's very very important...}
  \alert{... you can "alert" in an "alertblock"}\\
  Ewww, nasty, heh?
  \end{alertblock}
\end{frame}

\newcommand{\putlink}[1]{%
   \pgfsetlinewidth{1.4pt}%
   \pgfsetendarrow{\pgfarrowtriangle{4pt}}%
   \pgfline{\pgfxy(1,1)}{\pgfxy(#1,1)}
}

\begin{frame}
  \frametitle{Overlay effects}
  \framesubtitle{Keep the suspense!}
  \begin{block}{Time bomb}
  \begin{enumerate}
    \item<2-> Two more to go
    \item<3-> One more to go
    \item<4-> Last chance...
    \item<5-> BOOM!
  \end{enumerate}
  \end{block}
  \begin{block}{"Animation"}<6->
    \begin{pgfpicture}{0cm}{0cm}{7cm}{2cm}
    \only<1-6>{
      \putlink{2}
    }
    \only<7>{
      \putlink{4}
    }
    \only<8>{
      \putlink{6}
    }
    \only<9>{
      \putlink{8}
    }
    \only<10>{
      \putlink{10}
    }
    \end{pgfpicture}
  \end{block}
\end{frame}

\section*{}
\frame{
  \vfill
  \centering
  \highlighton{
  \usebeamerfont*{frametitle}And now?

  \usebeamerfont*{framesubtitle}Enter the secret section
  }
  \vfill
}
\begin{frame}
  \frametitle{Contributing to this beamer style}
  \framesubtitle{We want you !}

  \begin{block}{Why?}
  \begin{itemize}
    \item Beamer is hot!
    \item This style deserves to be improved
  \end{itemize}
  \end{block}

  \begin{block}{How?}
  \begin{itemize}
    \item Grab it
    \item Improve its LaTeX code
    \item Use you artistics skills
    \item Document it
    \item Help other people to use it
    \item Use it...
  \end{itemize}
  \end{block}
\end{frame}

\begin{frame}
  \frametitle{Resources}
  \framesubtitle{If you want to improve this style}
  \begin{thebibliography}{10}

  \beamertemplatearticlebibitems

  \bibitem{beamer-homepage}
    LaTeX Beamer
    \newblock {\tt http://latex-beamer.sourceforge.net/}

  \bibitem{kdeslides}
    KDE Presentations
    \newblock {\tt http://www.kde.org/kdeslides/}

  \end{thebibliography}
\end{frame}

\frame{
  \vspace{2cm}
  {\huge Questions ?}

}

\end{document}